\documentclass[../../../uem_tok_q.tex]{subfiles}

\begin{document}
    座標平面上に2点$\mathrm{O}(0,\;0)$，$\mathrm{A}(0,\;1)$をとる．%
    $x$軸上の2点$\mathrm{P}(p,\;0)$，$\mathrm{Q}(q,\;0)$が，%
    次の条件\ref{uem_tok_2024_f_hum_03_q:1}，\ref{uem_tok_2024_f_hum_03_q:2}をともに満たすとする．
    
    \begin{enumerate}[label=(\roman*), align=right,
        topsep=3pt, itemsep=3pt, itemindent=1\zw, leftmargin=2\zw, labelsep=.6\zw]
            \item $0<p<1$\;\;かつ\;\;$p<q$ \label{uem_tok_2024_f_hum_03_q:1}
            \item 線分APの中点をMとするとき，$\angle\mathrm{OAP}=\angle\mathrm{PMQ}$%
                \label{uem_tok_2024_f_hum_03_q:2}
    \end{enumerate}
    \begin{enumerate}
        \item $q$ を $p$ を用いて表せ．
        \item $q=\myfraca{1}{3}$ となる $p$ の値を求めよ．
        \item $\triangle\mathrm{OAP}$の面積を $S$，%
            $\triangle\mathrm{PMQ}$の面積を $T$ とする．$S>T$ となる $p$ の範囲を求めよ．
    \end{enumerate}
\end{document}
